Public evaluation assets for renewable-curtailment forecasting remain limited because studies commonly use private single-system data, task-specific targets, and incomparable protocols. We construct a reproducible curtailment-risk benchmark from the open Reliability Test System Grid Modernization Lab Consortium data. A fixed reference operating policy with a 70\% system non-synchronous-penetration-type cap converts 8760 hours of day-ahead load and renewable generation into a method-independent risk proxy, while an onset slice evaluates transitions from quiet to material-curtailment conditions. Under one leakage-free temporal split, training-window-calibrated warning thresholds, and ten seeded runs for stochastic methods, we compare DSTAR-GRU--a gated-recurrent operating-state retrieval framework--with eight baselines and five mechanism controls. Retrieval improves 1-hour mean absolute error relative to the prespecified retrieval-removal and retrieval-degradation controls, but reduces 24-hour onset-warning performance relative to no-retrieval and raw-feature-retrieval variants. Persistence has the lowest 1-hour mean absolute error. At 24 hours, SmallBank has the lowest mean absolute error but zero event F1; raw-feature kNN is the strongest non-degenerate mean-absolute-error reference, and ridge regression leads onset detection. A fixed-cap transport check on NREL-118 produces no positive target hours, indicating that the policy-derived task requires system-specific calibration. The benchmark therefore reveals a horizon-dependent role for retrieval and separates overall accuracy from operationally relevant onset warning without implying general forecasting superiority.
